\documentclass[./eval]{subfiles}
\usepackage[first=1, last=20]{lcg}
\renewcommand{\random}{\rand\arabic{rand}}

\begin{document}

\makebox[\textwidth]{  Esc.$N^{\circ}$:4-124 Reynaldo Merín $4^{to}1^{ra}$  \includegraphics[scale=0.04]{images/logos/4124.png} Evaluación Electrotecnia II | Apellido y Nombre:\enspace\hrulefill}

\begin{multicols}{2}
\begin{questions}
\question[2] El vacío tiene una permeabilidad magnética $\mu_0 =  4\pi \times 10^{-7}T\cdot$m/A ¿Cuál es la permeabilidad magnética relativa $\mu_r$ de un material que tiene una permeabilidad magnética absoluta de $\mu =\random \times 10^{-5}$ T$\cdot$m/A?
\question[1] Por un solenoide de $N=\random00$ espiras y longitud de $l=0,\random0m$ circula una corriente eléctrica de $I=\random0A$. Si el núcleo es vacío, ¿Cuál es la inducción magnética del solenoide?
\question[2] 
    Una bobina cilíndrica de superficie $S = \random cm^2$ cm lleva una corriente de $I = \random$ A tiene una longitud $l=\random cm$. Calcular el campo magnético si $N=\random00$ y la fuerza magneto-motríz.
    
\question[2] Si la bobina del problema anterior tiene un núcleo de hierro $\mu_{r}=\random00$. Calcular la inducción $B$, y el flujo $\Phi$ en la bobina en Maxwell[Mx].
\question[3] La intensidad de campo magnético en el interior del solenoide es de $H=\random00\frac{Av}{m}$ y su longitud es de $l=\random0cm$. Calcular:
    \begin{enumerate}
        \item Intensidad de corriente $I$, en el conductor si el solenoide  tiene $N=\random00$ espiras.        
    \end{enumerate}

\end{questions}
\end{multicols}





\end{document}